%% 第 6 章 Hecke 性质 II（原书 p.230--p.235）。
%% 2026-11-21 校订时把 units/ch06_u01 内联展开为本文件；此后本章不再由 assemble.py 覆盖。

% ===== units/ch06_u01.tex =====
% ch06_u01 —— 原书 p.230–p.235（第 6 章全部）
\bdchap{ch:hecke2}{Hecke 性质 II}
\SourcePage{230}

\bdsec{6.1.}
% [Tutor Note: 原书 p.230 的 6.1 节只有编号，既无标题也无正文，底本如此，保留原样。]

\bdsec{6.2. 定理 8.1.6 的证明}

\bdp{6.2.1} \emph{引理.} 设 $V$ 是非零 $\overline{U}'$-模，$\mathfrak{g}\otimes O$ 在 $V$ 上的表示可积，且理想 $I\subset\mathfrak{Z}$ 零化 $V$。则 $V$ 有非零的 $\mathfrak{g}\otimes O$-不变向量。

\emph{证明.} 记 $\mathfrak{m}$ 为 $O$ 的极大理想。态射 $G(O)\to G(O/\mathfrak{m})$ 的核是投射幂幺的，其李代数为 $\mathfrak{g}\otimes\mathfrak{m}$，故 $V^{\mathfrak{g}\otimes\mathfrak{m}}\neq 0$。考虑 Sugawara 元素 $\mathcal{L}_0\in I$（见 3.6.15、3.6.16）。由 (85) 一眼可见，$2\mathcal{L}_0$ 在 $V^{\mathfrak{g}\otimes\mathfrak{m}}$ 上的作用相当于 $\mathfrak{g}$ 的 Casimir 元。另一方面，因 $\mathcal{L}_0\in I$，有 $\mathcal{L}_0V=0$。所以 $\mathfrak{g}$ 在 $V^{\mathfrak{g}\otimes\mathfrak{m}}$ 上的作用平凡，且 $V^{\mathfrak{g}\otimes O}=V^{\mathfrak{g}\otimes\mathfrak{m}}\neq 0$。\hfill$\Box$

\bdp{6.2.2} \emph{引理.} 设 $N$ 是 $\mathfrak{z}_{\mathfrak{g}}(O)$-模，并带有李代数胚 $I/I^2$ 的作用。假设 $L_0\in\operatorname{Der}O\subset I/I^2$ 在 $N$ 上的作用可对角化，且对每个 $c\in\mathbb{C}$，其谱与 $c+\mathbb{Z}$ 的交有下界。则 $N$ 是自由 $\mathfrak{z}_{\mathfrak{g}}(O)$-模。

\emph{证明.} 用 (80)、(81) 与 3.6.17，可把 $\mathfrak{z}_{\mathfrak{g}}(O)$ 换成 $A_{{}^{L}\mathfrak{g}}(O)$，把 $I/I^2$ 换成 $\mathfrak{a}_{{}^{L}\mathfrak{g}}$。按定义，$\mathfrak{a}_{{}^{L}\mathfrak{g}}$ 是 $\mathfrak{F}^0_G$ 的无穷小对称构成的代数胚。在 3.5.6 中我们描述了 $\mathfrak{F}^0_G$ 的一个平凡化。相应的分裂 $\operatorname{Der}A_{{}^{L}\mathfrak{g}}(O)\to\mathfrak{a}_{{}^{L}\mathfrak{g}}$ 是 $\operatorname{Der}^0O$-等变的（见 (69) 与 (70)；要点在于，这些公式的右端作为 $\operatorname{Spec}A_{{}^{L}\mathfrak{g}}(O)$ 上的函数为常值）。于是 $N$ 成为 $\operatorname{Der}A_{{}^{L}\mathfrak{g}}(O)$ 上的模，且映射 $\operatorname{Der}A_{{}^{L}\mathfrak{g}}(O)\to\operatorname{End}N$ 是 $\operatorname{Der}^0O$-等变的。按 3.5.6，$A_{{}^{L}\mathfrak{g}}(O)$ 是 $u_{jk}$ 的多项式环，$1\le j\le r$，$0\le k<\infty$，且 $L_0u_{jk}=(d_j+k)u_{jk}$，其中 $d_j>0$。因此 $N$ 是由 $u_{jk}$ 与 $\partial/\partial u_{jk}$ 生成的代数上的 $L_0$-分次模，其中 $\deg(\partial/\partial u_{jk})=-\deg u_{jk}=-(d_j+k)\to-\infty$（$k\to\infty$）。于是 $N$ 的每个元素都被几乎所有的 $\partial/\partial u_{jk}$ 零化，也被 $\partial/\partial u_{jk}$ 的所有充分高次单项式零化。众所周知（例如见 [Kac90] 的引理 9.13，或 [Kac97] 的定理 3.5），
\SourcePage{231}
在此情形 $N=A_{{}^{L}\mathfrak{g}}(O)\otimes N_0$，其中 $N_0$ 由满足 $\partial/\partial u_{jk}\,n=0$（对一切 $j$、$k$）的 $n\in N$ 构成。\hfill$\Box$

\bdp{6.2.3} 下面证明定理 8.1.6。按 5.6.8，可把引理 6.2.2 用于 $N:=V^{\mathfrak{g}\otimes O}$。故 $N=\mathfrak{z}_{\mathfrak{g}}(O)\otimes W$，其中 $W$ 是某个向量空间。我们将证明自然的 $\overline{U}'$-模态射 $f:\mathit{Vac}'\otimes W=\mathit{Vac}'\otimes_{\mathfrak{z}_{\mathfrak{g}}(O)}N\to V$ 是同构。有 $(\operatorname{Ker}f)^{\mathfrak{g}\otimes O}=\operatorname{Ker}f\cap N=0$，故由 6.2.1 得 $\operatorname{Ker}f=0$。设 $\operatorname{Coker}f\neq 0$。则按 6.2.1，$\operatorname{Coker}f$ 中存在非零的 $\mathfrak{g}\otimes O$-不变元素，即非零的 $\overline{U}'$-模态射 $\mathit{Vac}'\to\operatorname{Coker}f$。它诱导扩张 $0\to\mathit{Vac}'\otimes W\to P\to\mathit{Vac}'\to 0$，该扩张不分裂（分裂 $\mathit{Vac}'\to P$ 与自然态射 $P\to V$ 的复合，会给出 $V$ 中不属于 $N$ 的 $\mathfrak{g}\otimes O$-不变向量）。于是只剩证明下述陈述。

\bdp{6.2.4} \emph{命题.} 离散 $\overline{U}'$-模的任一扩张 $0\to\mathit{Vac}'\otimes W\to P\to\mathit{Vac}'\to 0$，只要 $IP=0$，就分裂（这里 $W$ 是向量空间）。

\emph{证明.} 设 $p\in P$ 属于来自 $\mathit{Vac}'$ 的真空向量的原像。则 $(\mathfrak{g}\otimes O)\cdot p\subset\mathit{Vac}'\otimes W$。事实上，存在有限维子空间 $W_1\subset W$，使 $(\mathfrak{g}\otimes O)\cdot p\subset\mathit{Vac}'\otimes W_1$，故可设 $\dim W<\infty$。此外，由于函子 Ext 可加，可设 $W=\mathbb{C}$。

仍取上述 $p$。由 $\varphi(a)=ap$ 定义 $\varphi:\mathfrak{g}\otimes O\to\mathit{Vac}'$，于是 $\varphi$ 是 1-上闭链，且 $\operatorname{Ker}\varphi$ 为开。须证 $\varphi$ 是上边缘。有 $\overline{U}'$ 的标准滤链 $\overline{U}'_k$。$\mathit{Vac}'$ 上诱导的滤链 $\mathit{Vac}'_k$ 是 $(\mathfrak{g}\otimes O)$-不变的，因为真空向量被 $\mathfrak{g}\otimes O$ 零化。故 $\mathfrak{g}\otimes O$ 作用在 $\operatorname{gr}\mathit{Vac}'$ 上。存在 $k$，使得 $\operatorname{Im}\varphi\subset\mathit{Vac}'_k$。记 $\psi$ 为 $\varphi:\mathfrak{g}\otimes O\to\mathit{Vac}'_k$ 与 $\mathit{Vac}'_k\to\mathit{Vac}'_k/\mathit{Vac}'_{k-1}\subset\operatorname{gr}\mathit{Vac}'$ 的复合。于是 $\psi:\mathfrak{g}\otimes O\to\operatorname{gr}\mathit{Vac}'$ 是 1-上闭链，且只需证 $\psi$ 是上边缘（之后可用归纳法推进）。

\SourcePage{232}

记 $\mathit{Vac}^{cl}$ 为 $\mathfrak{g}^*\otimes\omega_O$ 上多项式构成的空间（按定义，$\mathfrak{g}^*\otimes\omega_O$ 上的多项式是函数 $\mathfrak{g}^*\otimes\omega_O\to\mathbb{C}$，它来自某个 $n$ 对应的向量空间 $\mathfrak{g}^*\otimes(\omega_O/\mathfrak{m}^n\omega_O)$ 上的多项式）。按 2.4.1，有典范的 $\mathfrak{g}\otimes O$-等变等同 $\operatorname{gr}\mathit{Vac}'=\operatorname{Sym}(\mathfrak{g}\otimes K/\mathfrak{g}\otimes O)=\mathit{Vac}^{cl}$（$\mathfrak{g}\otimes O$ 在 $\mathit{Vac}^{cl}$ 上的作用，由 $\mathfrak{g}\otimes O$ 在 $\mathfrak{g}^*\otimes\omega_O$ 上的自然作用诱导）。于是可把 $\psi$ 看作 1-上闭链 $\mathfrak{g}\otimes O\to\mathit{Vac}^{cl}$。由下式定义 $\beta_\psi:(\mathfrak{g}\otimes O)\times(\mathfrak{g}^*\otimes\omega_O)\to\mathbb{C}$：
\begin{equation}\tag{293}
\beta_\psi(a,\eta):=(\psi(a))(\eta).
\end{equation}

若 $\eta$ 在 $\mathfrak{g}^*\otimes(\omega_O/\mathfrak{m}\omega_O)$ 中的像正则，则称 $\eta\in\mathfrak{g}^*\otimes\omega_O$ 为\emph{正则}的。

\emph{引理.} 若 $\eta\in\mathfrak{g}^*\otimes\omega_O$ 正则，且 $\mathfrak{c}(\eta)$ 是 $\eta$ 在 $\mathfrak{g}\otimes O$ 中的稳定化子，则
\begin{equation}\tag{294}
\beta_\psi(a,\eta)=0,\qquad a\in\mathfrak{c}(\eta).
\end{equation}

\emph{证明.} 我们将用到 $IP=0$。设 $F\in\operatorname{Ker}(\mathfrak{Z}^{cl}\to\mathfrak{z}^{cl}_{\mathfrak{g}}(O))$，即 $F$ 是 $\mathfrak{g}^*\otimes\omega_K$ 上的 $(\mathfrak{g}\otimes K)$-不变多项式函数，且其在 $\mathfrak{g}^*\otimes\omega_O$ 上的限制为零（见 2.9.8）。设 $F$ 是 $r$ 次齐次的。按 3.7.8，$F$ 是某个 $z\in\mathfrak{Z}_r$ 的象征。由于 $F$ 在 $\mathfrak{z}^{cl}_{\mathfrak{g}}(O)$ 中的像为零，$z$ 在 $\mathfrak{z}_{\mathfrak{g}}(O)$ 中的像属于滤链的第 $(r-1)$ 项，故按 2.9.5，它来自某个 $z'\in\mathfrak{Z}_{r-1}$。把 $z$ 换成 $z-z'$，可设 $z\in I\cap\mathfrak{Z}_r$。

由于 $I\subset\overline{U}'\cdot(\mathfrak{g}\otimes O)$，可把 $z$ 写成
\begin{equation}\tag{295}
z=\sum_{i=1}^{\infty}u_ia_i,\qquad a_i\in\mathfrak{g}\otimes O,\quad u_i\in\overline{U}',\quad a_i\to 0\ (i\to\infty).
\end{equation}
由 Poincaré--Birkhoff--Witt 定理，可选取分解 (295)，使得对一切 $i$ 有 $u_i\in\overline{U}'_{r-1}$。把等式 $zp=0$ 改写为
\begin{equation}\tag{296}
\sum_i u_i\varphi(a_i)=0.
\end{equation}

\SourcePage{233}

记 $\tilde u_i$ 为 $u_i$ 在 $\overline{U}'_{r-1}/\overline{U}'_{r-2}$ 中的像。由 (295) 与 (296) 推出
\begin{equation}\tag{297}
F=\sum_i\tilde u_ia_i,
\end{equation}
\begin{equation}\tag{298}
\sum_i\tilde u_i\psi(a_i)=0
\end{equation}
其中 $a_i\in\mathfrak{g}\otimes O$ 被看作 $\mathfrak{g}^*\otimes\omega_K$ 上的线性函数，而 $\tilde u_i$ 是 $\tilde u_i$ 在 $\mathfrak{g}^*\otimes\omega_O$ 上的限制。
% [Tutor Note: 原书 p.233 此句两处均印作 $\tilde u_i$，第二个按文意当作 $u_i$；底本如此，保留原样。]

记 $dF$ 为 $F$ 的微分在 $\mathfrak{g}^*\otimes\omega_O$ 上的限制。由于 $F$ 在 $\mathfrak{g}^*\otimes\omega_O$ 上为零，有 $dF\in\mathit{Vac}^{cl}\widehat{\otimes}(\mathfrak{g}\otimes O)$，其中 $\widehat{\otimes}$ 是完备化张量积。按 (297)，$dF=\sum_i\tilde u_i\otimes a_i$，故可把 (298) 改写为
\begin{equation}\tag{299}
\mu(dF)=0
\end{equation}
其中 $\mu$ 是 $\operatorname{id}\otimes\psi:\mathit{Vac}^{cl}\widehat{\otimes}(\mathfrak{g}\otimes O)\to\mathit{Vac}^{cl}\otimes\mathit{Vac}^{cl}$ 与乘法映射 $\mathit{Vac}^{cl}\otimes\mathit{Vac}^{cl}\to\mathit{Vac}^{cl}$ 的复合。

现在令
\begin{equation}\tag{300}
F(\eta)=\operatorname{Res}f(\eta)\nu,\qquad \nu\in\omega_O^{\otimes(1-r)}
\end{equation}
其中 $f$ 是 $\mathfrak{g}^*$ 上 $r$ 次齐次不变多项式。此时 (299) 可改写为
\begin{equation}\tag{301}
\beta_\psi(A_f(\eta)\nu,\eta)=0
\end{equation}
其中 $\beta_\psi$ 由 (293) 定义，$A_f$ 是把 $f$ 看作多项式映射 $\mathfrak{g}^*\to\mathfrak{g}$ 时的微分（于是 $A_f(\eta)\in\mathfrak{g}\otimes\omega_O^{\otimes(r-1)}$，$A_f(\eta)\nu\in\mathfrak{g}\otimes O$）。由于 $f$ 不变，$A_f(l)$ 属于 $l\in\mathfrak{g}^*$ 的稳定化子；且若 $l$ 正则，则所有不变 $f$ 对应的 $A_f(l)$ 生成该稳定化子。于是引理由 (301) 推出。\hfill$\Box$

为证明命题，只剩说明：任一核为开的 1-上闭链 $\psi:\mathfrak{g}\otimes O\to\mathit{Vac}^{cl}$，只要函数 (293) 满足 (294)，就是上边缘。

\SourcePage{234}

\emph{引理.} 设 $K$ 是连通仿射代数群，$\operatorname{Hom}(K,\mathbb{G}_m)=0$，$W$ 是 $K$-模，$\psi$ 是 1-上闭链 $\operatorname{Lie}K\to W$。则 $\psi$ 来自唯一的 1-上闭链 $\Psi:K\to W$。

\emph{证明.} $\Psi$ 的唯一性显然。存在性的证明归结为 $K$ 幂幺的情形（把 $K$ 表示成半单子群 $K_{ss}$ 与幂幺正规子群的半直积；然后注意 $\psi$ 在 $\operatorname{Lie}K_{ss}$ 上的限制是上边缘，并归结为该限制为零的情形）。记 $\tilde K$ 为 $K$ 与 $W$ 的半直积。1-上闭链 $K\to W$ 等同于态射 $K\to\tilde K$，使得复合 $K\to\tilde K\to K$ 等于 $\operatorname{id}$。1-上闭链 $\operatorname{Lie}K\to W$ 有类似的解释。于是可用下述事实：函子 $\operatorname{Lie}:\{\text{幂幺群}\}\to\{\text{幂零李代数}\}$ 是等价。\hfill$\Box$

于是我们的 1-上闭链 $\psi:\mathfrak{g}\otimes O\to\mathit{Vac}^{cl}$ 来自 1-上闭链 $\Psi:G(O)\to\mathit{Vac}^{cl}$，这里 $G(O)$ 被看作群概形。由 $B_\Psi(g,\eta)=(\Psi(g))(\eta)$ 定义 $B_\Psi:G(O)\times(\mathfrak{g}^*\otimes\omega_O)\to\mathbb{C}$。

\emph{引理.} 若 $\eta\in\mathfrak{g}^*\otimes\omega_O$ 正则，且 $C(\eta)$ 是 $\eta$ 在 $G(O)$ 中的稳定化子，则
\begin{equation}\tag{302}
B_\Psi(g,\eta)=0,\qquad g\in C(\eta).
\end{equation}

\emph{证明.} 固定 $\eta$ 后，映射 $g\mapsto B_\Psi(g,\eta)$ 是群概形态射 $f:C(\eta)\to\mathbb{G}_a$。按 (294)，$f$ 的微分等于 $0$。故 $f=0$（即便 $C(\eta)$ 不连通，也有 $\operatorname{Hom}(\pi_0(C(\eta)),\mathbb{G}_a)=0$，因为 $\pi_0(C(\eta))$ 有限；但事实上，若取 $G$ 为伴随群（这不失一般性），则 $C(\eta)$ 连通）。\hfill$\Box$

$\Psi$ 是上闭链这一事实意味着
\begin{equation}\tag{303}
B_\Psi(g_1g_2,\eta)=B_\Psi(g_1,\eta)+B_\Psi(g_2,g_1^{-1}\eta g_1).
\end{equation}
须证 $B_\Psi$ 是上边缘，即
\begin{equation}\tag{304}
B_\Psi(g,\eta)=f(g^{-1}\eta g)-f(\eta)
\end{equation}
\SourcePage{235}
其中 $f$ 是某个多项式函数 $\mathfrak{g}^*\otimes\omega_O\to\mathbb{C}$。记 $\mathfrak{g}^*_{reg}$ 为 $\mathfrak{g}^*$ 的正则元构成之集，$(\mathfrak{g}^*\otimes\omega_O)_{reg}$ 为 $\mathfrak{g}^*\otimes\omega_O$ 的正则元构成之集（即 $\mathfrak{g}^*_{reg}$ 在 $\mathfrak{g}^*\otimes\omega_O$ 中的原像）。由于 $\operatorname{codim}(\mathfrak{g}^*\setminus\mathfrak{g}^*_{reg})>1$，只需在 $(\mathfrak{g}^*\otimes\omega_O)_{reg}$ 上构造正则函数 $f$。

设 $\mathcal{C}$ 的含义同 2.2.1。态射 $\mathfrak{g}^*_{reg}\to\mathcal{C}$ 光滑且满，$G$ 在其纤维上传递地作用；Kostant 在 [Ko63] 中构造了子概形 $\operatorname{Kos}\subset\mathfrak{g}^*_{reg}$，使得 $\operatorname{Kos}\to\mathcal{C}$ 是同构。若用 $\mathfrak{g}$ 上的不变纯量积把 $\mathfrak{g}^*$ 与 $\mathfrak{g}$ 等同，则 $\operatorname{Kos}=i\left(\begin{pmatrix}0&0\\1&0\end{pmatrix}\right)+V$，其中 $i$ 与 $V$ 的含义同 3.1.9。由 $\operatorname{Kos}_O:=i\left(\begin{pmatrix}0&0\\1&0\end{pmatrix}\right)\cdot dt+V\otimes\omega_O$ 定义 $\operatorname{Kos}_O\subset\mathfrak{g}^*\otimes\omega_O$。

方程 (304) 有唯一解 $f$，它在 $\operatorname{Kos}_O$ 上为零。$f$ 在 $(\mathfrak{g}^*\otimes\omega_O)_{reg}$ 上的限制由下式定义：
\begin{equation}\tag{305}
f(g^{-1}\eta g)=B_\Psi(g,\eta),\qquad \eta\in\operatorname{Kos},\ g\in G(O).
\end{equation}

这里 $f$ 良定义，因为（由 (302) 与 (303) 可得）对 $\eta\in(\mathfrak{g}^*\otimes\omega_O)_{reg}$ 与 $g_1\in C(\eta)$ 有 $B_\Psi(g_1g,\eta)=B_\Psi(g,\eta)$。于是 (303) 表明，由 (305) 定义的函数 $f$ 满足 (304)。\hfill$\Box$

\emph{注.} 证明末尾我们用到了 Kostant 给出的纤维化 $(\mathfrak{g}^*\otimes\omega_O)_{reg}\to\operatorname{Hitch}_{\mathfrak{g}}(O)$ 的整体截面（$\operatorname{Hitch}_{\mathfrak{g}}(O)$ 的定义见 2.4.1）。也可改用\emph{局部}截面与等式 $H^1(\operatorname{Hitch}_{\mathfrak{g}}(O),\mathcal{O})=0$，后者显然，因为 $\operatorname{Hitch}_{\mathfrak{g}}(O)$ 是仿射的。

\bdp{6.2.5} 命题 6.2.4 似乎与 [F91] 有关（例如见 [F91] 第 97、98 页下半部分的命题）。或许修改 [F91] 的方法，可得到命题 6.2.4 乃至更多。
